\section{Arithmétique}

\begin{UPSTIexercice}{Nombre premier}
	\UPSTIetape Écrire un programme \texttt{premier.c} qui demande un entier \texttt{n} (\texttt{n >= 2}) et affiche s'il est premier, en testant ses diviseurs à l'aide d'une boucle. Réfléchir à la condition d'arrêt pour ne pas tester plus de diviseurs que nécessaire.
	Par exemple :
	\begin{lstlisting}[language=bash,style=console]
Entrez un entier : 17
17 est premier
\end{lstlisting}
\end{UPSTIexercice}

\begin{UPSTIcorrectionP}{Nombre premier}
	\begin{lstlisting}[language=c]
#include <stdio.h>

int main(void) {
    int n;
    printf("Entrez un entier : ");
    scanf("%d", &n);

    int est_premier = 1;
    for (int d = 2; d * d <= n; d++) {
        if (n % d == 0) {
            est_premier = 0;
            break;
        }
    }

    if (est_premier) {
        printf("%d est premier\n", n);
    } else {
        printf("%d n'est pas premier\n", n);
    }

    return 0;
}
	\end{lstlisting}
\end{UPSTIcorrectionP}

\begin{UPSTIexercice}{PGCD par l'algorithme d'Euclide}
	\UPSTIetape Écrire un programme \texttt{pgcd.c} qui demande deux entiers strictement positifs \texttt{a} et \texttt{b}, et calcule leur PGCD (plus grand commun diviseur) avec une boucle \texttt{while}, selon l'algorithme d'Euclide : tant que le reste n'est pas nul, remplacer le plus grand nombre par le reste de la division de l'un par l'autre.
	Par exemple :
	\begin{lstlisting}[language=bash,style=console]
Entrez a : 48
Entrez b : 18
PGCD(48, 18) = 6
\end{lstlisting}
\end{UPSTIexercice}

\begin{UPSTIcorrectionP}{PGCD par l'algorithme d'Euclide}
	\begin{lstlisting}[language=c]
#include <stdio.h>

int main(void) {
    int a, b;
    printf("Entrez a : ");
    scanf("%d", &a);
    printf("Entrez b : ");
    scanf("%d", &b);

    int x = a, y = b;
    while (y != 0) {
        int reste = x % y;
        x = y;
        y = reste;
    }

    printf("PGCD(%d, %d) = %d\n", a, b, x);

    return 0;
}
	\end{lstlisting}
\end{UPSTIcorrectionP}

\begin{UPSTIexercice}{Suite de Syracuse}
	\UPSTIetape Écrire un programme \texttt{syracuse.c} qui demande un entier strictement positif \texttt{n} et applique la règle de Syracuse jusqu'à atteindre \texttt{1} : si le nombre courant est pair, le diviser par 2 ; sinon, le remplacer par \texttt{3n + 1}. Le programme affichera le nombre d'étapes effectuées et la plus grande valeur atteinte durant le parcours.
	Par exemple :
	\begin{lstlisting}[language=bash,style=console]
Entrez un entier : 6
6 -> 3 -> 10 -> 5 -> 16 -> 8 -> 4 -> 2 -> 1
8 etapes, valeur maximale atteinte : 16
\end{lstlisting}
\end{UPSTIexercice}

\begin{UPSTIcorrectionP}{Suite de Syracuse}
	\begin{lstlisting}[language=c]
#include <stdio.h>

int main(void) {
    int n;
    printf("Entrez un entier : ");
    scanf("%d", &n);

    int valeur = n;
    int etapes = 0;
    int maximum = valeur;

    printf("%d", valeur);
    while (valeur != 1) {
        if (valeur % 2 == 0) {
            valeur /= 2;
        } else {
            valeur = 3 * valeur + 1;
        }
        etapes++;
        if (valeur > maximum) {
            maximum = valeur;
        }
        printf(" -> %d", valeur);
    }
    printf("\n%d etapes, valeur maximale atteinte : %d\n", etapes, maximum);

    return 0;
}
	\end{lstlisting}
\end{UPSTIcorrectionP}
